\section*{Resumen}

El modelo de Ising es, desde casi su concepción, uno de los sistemas más importantes de estudio en la mecánica estadística. Esto es debido a que desde su increíble simplicidad, es capaz de presentar una amplia gama de fenómenos distintos y relevantes, como la transición de fase ferromagnética-paramagnética de un material 2D, la divergencia de la longitud de correlación en  la temperatura crítica $T_c$ o el ciclo de histéresis magnética; y propiedades como la clase de universalidad y la existencia de exponentes críticos correspondientes a esta.

En este trabajo se estudia el modelo bidimensional mediante el método de Montecarlo, usando y comparando tres algoritmos distintos: Metropolis, Glauber y Wolff. La motivación detrás de implementar más de un algoritmo es que cada uno tiene sus puntos negativos y positivos, por lo que se hace un estudio comparativo para diferenciar bajo qué circunstancias resultan mejores unos u otros.

Las simulaciones se realizan sobre una red cuadrada de lados $L\in \{16,32,64,128\} $ con condiciones de contorno periódicas, tomando el acoplamiento magnético como $J=1$ y la temperatura en unidades de $J/k_B$. Para cada tamaño, temperatura y red se recogen suficientes muestras tras un período de termalización como para intentar calcular con precisión las magnitudes buscadas. Entre los observables calculados se encuentran la magnetización media por partícula $\langle |m| \rangle$, la capacidad calorífica $C$, la susceptibilidad magnética $\chi$, el cumulante de cuarto orden de Binder $U_4$ y la longitud de correlación de segundo momento $\xi_L$. 

Para el estudio de los exponentes críticos se hace uso de la técnica del escalado de tamaño finito (\textit{Finite Size Scaling}, FSS): si los datos se reescalan adecuadamente usando las potencias de $L$ usando los exponentes críticos, la teoría dice que las curvas de distintos tamaños deberían colapsar unas sobre las otras. Se ajustan los exponentes para encontrar el colapso de los datos de manera visual, y con ello los valores aproximados de los exponentes críticos obtenidos con nuestra simulación. Este colapso se comprueba para la susceptibilidad magnética y la longitud de correlación.

El cumulante de Binder ya es una magnitud escalada, y tiene la propiedad de que las curvas para distintos tamaños se cortan aproximadamente en la temperatura crítica, lo que permite obtener un resultado sin la necesidad de ajustar exponentes. Sin embargo, aún así implementamos un segundo método de cálculo de la temperatura crítica usando el cumulante de Binder, para intentar eliminar los posibles efectos de tamaño finito.

El otro aspecto relevante del trabajo es la comparación entre los tres algoritmos utilizados. El tiempo de autocorrelación integrado $\tau_{\text{int}}$, calculado mediante el estimador de ventana de Madras--Sokal \citep{Madras1988}, muestra que Metropolis y Glauber sufren un fenómeno conocido como \textit{critical slowing down}: se necesitan muchas más medidas independientes cerca de $T_c$ para obtener un resultado significativo, algo característico de las dinámicas locales. Cabe destacar que Glauber sufre más este fenómeno que Metropolis, lo que nos permite ver las diferencias entre los mecanismos de actualización de ambas. Wolff es lo que se conoce como un algoritmo de \textit{cúmulo}, que elimina este fenómeno: $\tau_{\text{int}}$ de Wolff es muy bajo y no tiene estructura en todo el rango de temperaturas analizado.

Los resultados obtenidos son satisfactorios al compararlos con lo predicho por la teoría, desde los exponentes y la temperatura crítica de Onsager a la distinción entre diferentes algoritmos en magnitudes como el tiempo o la longitud de correlación.
\enlargethispage{\baselineskip}

\textbf{Palabras clave:} modelo de Ising, Monte Carlo, Metropolis, Glauber, Wolff, transición de fase, escalado de tamaño finito, exponentes críticos, temperatura crítica, \textit{critical slowing down}.

\section*{Abstract}
The Ising model has been, almost since its conception, one of the most important systems of study in statistical mechanics. This is due to the fact that, from its incredible simplicity, it is capable of presenting a wide range of different and relevant phenomena, such as the ferromagnetic-paramagnetic phase transition of a 2D material, the divergence of the correlation length at the critical temperature $T_c$, the universality class and the existence of critical exponents, or the magnetic hysteresis cycle.

In this work the two-dimensional model is studied using the Monte Carlo method, implementing and comparing three different algorithms: Metropolis, Glauber and Wolff. The motivation behind implementing more than one algorithm is that each one has its own drawbacks and advantages, so a comparative study is carried out to differentiate under which circumstances each one performs better.

The simulations are performed on square lattices with side lengths $L\in \{16,32,64,128\}$ with periodic boundary conditions, using the magnetic coupling $J=1$ and temperature in units of $J/k_B$. For each size, temperature and lattice, enough samples are collected after a thermalisation period in order to attempt to calculate the target quantities with sufficient precision. Among the computed observables are the mean magnetisation per spin $\langle |m| \rangle$, the heat capacity $C$, the magnetic susceptibility $\chi$, the fourth-order Binder cumulant $U_4$, and the second-moment correlation length $\xi_L$.

To study the critical exponents, the Finite Size Scaling (FSS) technique is employed: if the data are appropriately rescaled using powers of $L$ with the critical exponents, theory predicts that the curves for different system sizes should collapse onto one another. The exponents are fitted to find the data collapse visually, thereby obtaining approximate values of the critical exponents from our simulation. This collapse is verified for both the magnetic susceptibility and the correlation length.

The Binder cumulant is already a scaled quantity, and has the property that curves for different system sizes cross approximately at the critical temperature, which allows one to obtain a result without the need to fit any exponents. Nevertheless, we also implement a second method for computing the critical temperature using the Binder cumulant, in order to attempt to eliminate possible finite-size effects.

The other relevant aspect of the work is the comparison between the different algorithms used. The integrated autocorrelation time $\tau_{\text{int}}$, computed using the Madras--Sokal windowing estimator \citep{Madras1988}, shows that Metropolis and Glauber suffer from a phenomenon known as critical slowing down: many more independent measurements are needed near $T_c$ to obtain a meaningful result, which is characteristic of local dynamics. It is worth noting that Glauber suffers more from this phenomenon than Metropolis, which allows us to observe the differences between the update mechanisms of both. Wolff is what is known as a cluster algorithm, which eliminates this phenomenon: the $\tau_{\text{int}}$ of Wolff is very low and has no structure across the entire analysed temperature range.
The results obtained are satisfactory when compared with theoretical predictions, from Onsager's critical exponents and critical temperature to the distinction between different algorithms in quantities such as the autocorrelation time or the correlation length.

\textbf{Keywords:} Ising model, Monte Carlo, Metropolis, Glauber, Wolff, phase transition, finite size scaling, critical exponents, critical temperature, critical slowing down.